% Terminal apparatus of the concise study: the scope note and the References
% for the sources this companion actually uses. The revision timestamp and the
% rights colophon follow from the shared generation record.

\section{Appendix: Scope and Qualifications}

{\footnotesize
\textbf{Edition, formulary and occurrence.} This concise study treats the Mass
of the Twenty-ninth Sunday in Ordinary Time, Year~A, as the \work{Roman
Missal}, Third Edition, whose Latin base is the \work{Missale Romanum}, editio
typica tertia (2002), and no.~145 of the \work{Lectionary for Mass for Use in
the Dioceses of the United States of America}, which follows the \work{Ordo
lectionum Missae} of 1981, give it for the dioceses of the United States, in
the national calendar alone, for Sunday 18~October 2026, when the Feast of
Saint Luke is omitted (\work{Universal Norms} 60). Both Communion antiphons
are treated; a
claim resting on one holds only where it is chosen, and a claim resting on an
antiphon's wording does not hold where another approved chant is sung in its
place. Nothing here holds for a celebration that replaces this formulary with
another Mass.

\textbf{Relation to the full study.} Every claim here is drawn from the full
study of this same formulary, the Research PDF subtitled \textit{Caesar's coin
and God's image: a study of the proper in three interpretations}, which
develops the three interpretations at length, each with its own four senses;
nothing evidence-dependent is added here. The four rows on the
first page orient a reader and do not replace those senses. Each Father and
later writer is cited only for what he said of his own passage.

\textbf{Relations among the elements.} The liturgical books state four
relations: the first reading is chosen in relation to the Gospel, as the
\work{Ordo}'s titles show (\work{Praenotanda} 105--107); the psalm responds to
the first reading (\work{General Instruction} 61); the Alleluia verse greets
the Gospel (\work{General Instruction} 62); and the second reading
and the Gospel each belong to a semi-continuous course. The witnesses' own
joins are those named beside them (Augustine's Psalm~4 and 1~John 3:2,
Chrysostom's Romans 13 and Acts 17, Aquinas's Baruch, Bede's and Chrysostom's
context of Mark 10:45). Every other connection between
one element and another is the editor's synthesis. No witness cited comments
on this Mass or on any of its three orations, and none joins Isaiah 45 to
Matthew 22.

\textbf{Text, English and rights.} Scripture is quoted in the Douay--Rheims
version as revised by Bishop Challoner, a public-domain study translation of
the Clementine Vulgate, and not in the approved English proclaimed at Mass,
which is that of the Lectionary; the Missal's antiphons and orations are
identified by Latin incipit and described, and the approved English of neither
book is reproduced: only incipits, the \work{Ordo}'s titles and short Latin
phrases are quoted from them. English renderings of Latin sources are the editor's. The
other editions and translations named below are in the public domain, except
the Corpus Thomisticum texts and \work{Nostra aetate}, quoted in short
phrases, the \work{General Instruction}, cited only, and the introductions of
the \work{New American Bible, Revised Edition}, summarized.

\textbf{The questioners.} Several Fathers read for the full study pass, in
sentences near those quoted, from the men who questioned Jesus to the Jewish
people as a whole. None of those sentences is used. The agents are those the
Gospel names, the disciples of the Pharisees and the Herodians; hypocrisy is
treated as a vice that belongs to anyone; and the Second Vatican Council's
teaching that the Jews are not to be presented as rejected or accursed by God,
as if this followed from the Scriptures, is followed (\work{Nostra aetate} 4).

\textbf{Reception and its limits.} Augustine, Chrysostom, Justin Martyr,
Cyprian, Novatian, \work{Barnabas}, Cyril and Bellarmine were read in the
English translations named below; Jerome, Hilary, Bede and Aquinas on Matthew,
Isaiah and Philippians on page images of their printings; Ambrose and the
Corpus Thomisticum texts in electronic Latin. No witness was collated with a
modern critical edition. Augustine's
whited wall (\work{Enarrationes in Psalmos} 95.3) is used for its definition of
hypocrisy only, without the application he gives it in that section. Cyril,
Justin and Cyprian carry no interpretation; Cyprian, Novatian and
\work{Barnabas} witness only to the variant of Isaiah 45:1. Augustine's
``some'' shows only that the image reading was known to him.

\textbf{Chronology and records.} The rows on the second page are the generated
projection of the Scripture chronology corpus for this formulary, with its
relations, alternatives and unresolved states unchanged; the place phrases
rest on the sources recorded in \texttt{research/scope.md}. The dates come from articles of the \work{Catholic Encyclopedia}, the
introductions of the \work{New American Bible, Revised Edition} to Isaiah and
the Psalms, and the Pontifical Biblical Commission's answer of 1912 on Mark.
The search, its loci and its negative results are in
\texttt{research/scope.md}, the reasoning behind the three interpretations in
\texttt{research/interpretations.md}, and this companion's production record in
\texttt{research/production-review.md}.
\par}

\phantomsection
\section*{References}
\runninghead{References}

{\scriptsize
\subsection*{Liturgical books and Bibles}
\begin{itemize}\setlength{\itemsep}{0.15em}\setlength{\parskip}{0pt}
\item \work{Missale Romanum}, editio typica tertia (Typis Vaticanis, 2002), \latin{Dominica XXIX ``per annum''}; the \latin{Normae universales de anno liturgico et de calendario} 60, as printed in that Missal.
\item \work{The Roman Missal}, Third Edition, for Use in the Dioceses of the United States of America (2011).
\item \work{Ordo lectionum Missae}, editio typica altera (Libreria Editrice Vaticana, 1981): \work{Praenotanda} 105--107; no.~145, \latin{Dominica vigesima nona}, Year~A.
\item \work{Lectionary for Mass for Use in the Dioceses of the United States of America}, second typical edition, vol.~I, no.~145, as witnessed by the USCCB \work{Liturgical Calendar} for 2026 and page of readings for 18~October 2026.
\item \work{General Instruction of the Roman Missal}, United States edition of 2011 with the emendations of 2021, nos.~61--62.
\item The Holy Bible, Douay--Rheims version, revised by Bishop Richard Challoner (Project Gutenberg text): Gn 1:26; Ps 16:1--8; 32:15--19; 95:1--13; Is 45:1--6; Mt 21:45--46; 22:15--21; Mk 10:42--45; Acts 17:7; Rom 13:7; 1~Cor 15:49; Phil 2:15--16; 1~Thes 1:1--5; Jas 1:8.
\item \work{Biblia Sacra Vulgatae editionis} (Clementine Vulgate), eBible.org text: Ps 16:6--8; Is 45:1, 4; Mt 22:17--21; Phil 2:15.
\end{itemize}

\subsection*{Fathers, Doctors and saints}
\begin{itemize}\setlength{\itemsep}{0.15em}\setlength{\parskip}{0pt}
\item St Augustine, \work{Enarrationes in Psalmos} 4.8; 16.2, 8; 57.8; 95.2--3, 8--9; 115.13; 124.7--8; Nicene and Post-Nicene Fathers, 1st ser., vol.~8.
\item St Augustine, \work{In Iohannis evangelium tractatus} 40.9 and 41.2, Nicene and Post-Nicene Fathers, 1st ser., vol.~7; \work{De sermone Domini in monte} II.20.70, vol.~6.
\item St Ambrose, \work{Expositio evangelii secundum Lucam} IX.34--35 (Migne, PL 15), in an electronic Latin text of that printing.
\item St Hilary of Poitiers, \work{Commentarius in Matthaeum} 23.1--2 (Migne, PL 9, Paris, 1844, cols.~1044--1045).
\item St Jerome, \work{Commentarii in Matthaeum} III, on 22:15--21 (Migne, PL 26, Paris, 1845, cols.~162--163); \work{Commentarii in Isaiam} XII, on 45:1--7 (Migne, PL 24, Paris, 1865, cols.~457--458).
\item St John Chrysostom, \work{Homilies on Matthew} 6, 65.4 and 70.1--2, Nicene and Post-Nicene Fathers, 1st ser., vol.~10; \work{Homilies on Romans} 23, vol.~11; \work{Homilies on Philippians} 8 and \work{Homilies on First Thessalonians} 1, vol.~13.
\item St Bede, \work{In Marci evangelium expositio} III, on 10:42--45 (Migne, PL 92, Paris, 1862, col.~236).
\item St Thomas Aquinas, \work{Super Evangelium Matthaei} c.~22, in \work{Opera}, editio altera Veneta, vol.~3 (Venice, 1745), pp.~283--284; \work{Expositio super Isaiam ad litteram} c.~45, in \work{Opera omnia}, vol.~14 (Parma, 1863), p.~539; \work{Super epistolam ad Philippenses} c.~2, lect.~4, in \work{In omnes S.~Pauli Apostoli epistolas commentaria}, vol.~2 (Li\`ege: Dessain, 1857), p.~383.
\item St Thomas Aquinas, \work{Super I epistolam ad Thessalonicenses} c.~1, lect.~1, and \work{Postilla super Psalmos} 32 (Corpus Thomisticum text).
\item St Robert Bellarmine, \work{A Commentary on the Book of Psalms}, trans. J. O'Sullivan (1866, abridged), in the eCatholic2000 transcription: on Ps 32:18--19; 95:5, 10, 13.
\end{itemize}

\subsection*{Other witnesses, cited in support, and the Magisterium}
\begin{itemize}\setlength{\itemsep}{0.15em}\setlength{\parskip}{0pt}
\item Cyril of Alexandria, \work{Commentary on the Gospel of Saint Luke}, Sermon 135, trans. R. Payne Smith (Oxford, 1859).
\item St Justin Martyr, \work{First Apology} 17, and the \work{Epistle of Barnabas} 12, Ante-Nicene Fathers, vol.~1; St Cyprian, \work{Ad Quirinum} (\work{Testimonia}) I.21, and Novatian, \work{De Trinitate} 26, vol.~5.
\item Second Vatican Council, Declaration on the Relation of the Church to Non-Christian Religions \work{Nostra aetate} (28~October 1965), no.~4, Holy See website.
\end{itemize}

\subsection*{Chronology sources}
\begin{itemize}\setlength{\itemsep}{0.15em}\setlength{\parskip}{0pt}
\item \work{The Catholic Encyclopedia} (New York, 1908--1912): Gerard, ``General Chronology'', and Howlett, ``Biblical Chronology'' (vol.~3); Corbett, ``David'' (vol.~4); Souvay, ``Isaias'' (vol.~8); MacRory, ``Gospel of Saint Mark'' (vol.~9); Jacquier, ``Gospel of St.~Matthew'' (vol.~10); Prat, ``St.~Paul'' (vol.~11); Drum, ``Psalms'', and Vander Heeren, ``Epistle to the Philippians'' (vol.~12); Drum, ``Epistles to the Thessalonians'', Durand, ``The New Testament'', and Meistermann, ``Temple of Jerusalem'' (vol.~14); read in the New Advent text.
\item United States Conference of Catholic Bishops, \work{New American Bible, Revised Edition}, introductions to Isaiah and to the Psalms (summarized; protected).
\item Pontifical Biblical Commission, \latin{De auctore, de tempore compositionis et de historica veritate Evangeliorum secundum Marcum et secundum Lucam}, \work{Acta Apostolicae Sedis} 4 (1912), pp.~463--465, as the chronology corpus records it.
\end{itemize}
\par}
